% 级数（综述）
% license CCBYSA3
% type Wiki

本文根据 CC-BY-SA 协议转载翻译自维基百科\href{https://en.wikipedia.org/wiki/Series_(mathematics)}{相关文章}

在数学中，级数大致而言就是将无限多个项一个接一个地相加。[1]对级数的研究是微积分及其推广——数学分析——的重要组成部分。级数在数学的大多数领域都有应用，甚至在组合学中研究有限结构时，也会通过生成函数来使用级数。无限级数的数学性质使其在物理学、计算机科学、统计学和金融等其他定量学科中也得到了广泛应用。

在古希腊人看来，一个“潜在无限”的求和可以得到一个有限的结果，这被认为是悖论，其中最著名的例子就是芝诺悖论。[2][3]尽管如此，古希腊数学家们，包括阿基米德，仍然在实践中使用了无限级数，例如在抛物线求积问题中。[4][5]芝诺悖论中的数学问题在 17 世纪通过“极限”概念得到了化解，尤其是通过艾萨克·牛顿的早期微积分。[6]在19世纪，通过卡尔·弗里德里希·高斯和奥古斯丁-路易·柯西等人的工作，这一解决方案得到了进一步的严格化与改进。[7]他们回答了哪些和式是存在的问题（通过实数的完备性），以及级数项是否可以重新排列而不改变其和的问题（通过绝对收敛与条件收敛的概念）。

在现代术语中，任何一个有序的无限序列$(a_{1}, a_{2}, a_{3}, \ldots)$，无论这些项是数字、函数、矩阵，还是其他可以相加的对象，都定义了一个级数，即把这些$a_i$ 一个接一个地相加。为了强调项的数量是无限的，级数常常也被称为无限级数，以区别于有限级数，后者有时用来指有限和。级数通常写作：
$$
a_{1} + a_{2} + a_{3} + \cdots ,~
$$
或者使用大写希腊字母 Σ 的求和符号表示：[8]
$$
\sum_{i=1}^{\infty} a_i .~
$$
由于级数所表达的无限次加法无法在有限时间内逐一完成，但如果这些项及其有限部分和属于一个具有极限的集合，就可能为该级数赋予一个值，称为级数的和。这个值定义为：当 $n \to \infty$ 时，级数前$n$项的部分和的极限（若该极限存在）。[9][10][11] 这些有限和称为部分和。用求和符号表示：
$$
\sum_{i=1}^{\infty} a_i = \lim_{n \to \infty} \sum_{i=1}^{n} a_i ,~
$$
若极限存在。[9][10][11]当极限存在时，该级数称为收敛或可和，此时序列$(a_{1}, a_{2}, a_{3}, \ldots)$ 也是可和的；反之，如果极限不存在，则称该级数发散。[9][10][11]
表达式$\sum_{i=1}^{\infty} a_{i}$既表示级数——即无限地一个接一个相加项的隐含过程——又在该级数收敛时表示级数的和——即这一过程的明确极限。这是对类似约定的一种推广：例如$a+b$既表示“加法”这一过程，也表示其结果——$a$ 与 $b$ 的和。

通常，级数的项来自某个环，最常见的是实数域$\mathbb{R}$或复数域$\mathbb{C}$。在这种情况下，所有级数的集合本身也构成一个环，其中的加法是逐项相加，而乘法则由柯西乘积定义。[12][13][14]
\subsection{定义}
\subsubsection{级数}
一个级数，或更完整地说，无限级数，就是一个无限和。它通常表示为：[8][15][16]
$$
a_{0} + a_{1} + a_{2} + \cdots \quad \text{或} \quad a_{1} + a_{2} + a_{3} + \cdots ,~
$$
其中各项 $a_{k}$ 是一个数列（序列）的成员，这些成员可以是数、函数，或任何可以相加的对象。级数也常用大写希腊字母 Σ 的符号来表示：[8][16]
$$
\sum_{k=0}^{\infty} a_{k} \quad \text{或} \quad \sum_{k=1}^{\infty} a_{k}.~
$$
另外，一种常见的表示方法是写出前若干项，加上省略号，再写出一般项，最后再加省略号，其中一般项是$n$的函数形式的第$n$项：
$$
a_{0} + a_{1} + a_{2} + \cdots + a_{n} + \cdots 
\quad \text{或} \quad 
f(0) + f(1) + f(2) + \cdots + f(n) + \cdots .~
$$
例如，欧拉常数$e$可以通过以下级数定义：
$$
\sum_{n=0}^{\infty} \frac{1}{n!} 
= 1 + 1 + \frac{1}{2} + \frac{1}{6} + \cdots + \frac{1}{n!} + \cdots ,~
$$
其中$n!$表示前$n$个正整数的乘积，而按照约定$0! = 1$。[17][18][19]
\subsubsection{级数的部分和}
给定一个级数$s = \sum_{k=0}^{\infty} a_k $,它的第 $n$ 个**部分和**定义为：[9][10][11][16]
$$
s_{n} = \sum_{k=0}^{n} a_k = a_0 + a_1 + \cdots + a_n .~
$$
有些作者会直接把一个级数与其部分和序列等同。[9][11] 无论是部分和序列还是项的序列，都能完全刻画这个级数。而且，项的序列可以通过部分和序列的相邻差来恢复：
$$
a_{n} = s_{n} - s_{n-1}.~
$$
数列的部分求和是线性序列变换的一个例子，在计算机科学中也称为前缀和。而从部分和恢复原数列的逆变换就是有限差分，这也是另一种线性序列变换。

有时，级数的部分和具有更简单的闭合形式。例如：等差级数的部分和：
$$
s_{n} = \sum_{k=0}^{n} (a + kd) 
= a + (a+d) + (a+2d) + \cdots + (a+nd) 
= (n+1)\left(a + \tfrac{1}{2}nd\right) .~
$$
等比级数的部分和：[20][21][22]
$$
s_{n} = \sum_{k=0}^{n} ar^{k} 
= a + ar + ar^{2} + \cdots + ar^{n} 
= a \cdot \frac{1-r^{\,n+1}}{1-r}, \quad \text{当 } r \neq 1,~
$$
或者更简单地写成：
$$
s_{n} = a(n+1), \quad \text{当 } r = 1.~
$$
\subsubsection{级数的和}
\begin{figure}[ht]
\centering
\includegraphics[width=10cm]{./figures/59c0b57e291abb28.png}
\caption{三种几何级数及其从第 1 项到第 6 项部分和的示意图。虚线表示极限。} \label{fig_JiSu_1}
\end{figure}
严格来说，当一个级数的部分和序列存在极限时，就称这个级数收敛、是收敛的，或是可和的。当部分和序列的极限不存在时，就称该级数发散或是发散的。[23]当部分和的极限存在时，它被称为级数的和或级数的值：[9][10][11][16]
$$
\sum_{k=0}^{\infty} a_{k} 
= \lim_{n \to \infty} \sum_{k=0}^{n} a_{k} 
= \lim_{n \to \infty} s_{n}.~
$$
只有有限个非零项的级数总是收敛的。这样的级数在处理有限和时很有用，无需关心项数的多少。[24]

当和存在时，级数的和与它的第$n$个部分和之间的差：$s - s_{n} = \sum_{k=n+1}^{\infty} a_{k}$,称为该无限级数的第$n$阶截断误差。[25][26]

一个收敛级数的例子是几何级数：
$$
1 + \tfrac{1}{2} + \tfrac{1}{4} + \tfrac{1}{8} + \cdots + \tfrac{1}{2^{k}} + \cdots .~
$$
通过代数运算可以证明，它的第 $n$ 个部分和 $s_{n}$ 为：
$$
\sum_{k=0}^{n} \tfrac{1}{2^{k}} = 2 - \tfrac{1}{2^{n}} .~
$$
由于
$$
\lim_{n \to \infty} \left( 2 - \tfrac{1}{2^{n}} \right) = 2 ,~
$$
该级数收敛，并且其和为$2$，其第$n$阶截断误差为$\tfrac{1}{2^{n}}$。[20][21][22]

相比之下，几何级数
$$
\sum_{k=0}^{\infty} 2^{k}~
$$
在实数范围内是发散的。[20][21][22] 然而，在扩展实数轴上，它是收敛的，其极限为 $+\infty$，并且在每一步的截断误差也为 $+\infty$。[27]

当一个级数的部分和序列不易直接计算和检验其收敛性时，可以使用收敛性判别法来证明该级数是收敛还是发散。
\subsection{分组与重排项}
\subsubsection{分组}
在普通的有限求和中，由于加法的结合律，和式中的项可以任意分组或拆分，而不会改变求和的结果：
$$
a_{0} + a_{1} + a_{2} 
= a_{0} + (a_{1} + a_{2}) 
= (a_{0} + a_{1}) + a_{2}.~
$$
类似地，在一个级数中，任何有限次分组都不会改变其部分和的极限，因此不会改变级数的和。但是，如果在一个无限级数中进行无限次分组，那么分组后的部分和可能会有不同于原级数的极限，而不同的分组方式之间也可能得到不同的极限。例如：原和式$a_{0} + a_{1} + a_{2} + \cdots$的和不一定等于分组形式$a_{0} + (a_{1}+a_{2}) + (a_{3}+a_{4}) + \cdots $.

例如，格兰迪级数$1 - 1 + 1 - 1 + \cdots$的部分和序列在 1 与 0 之间来回交替，因此它不收敛。若将项成对分组，可得：$(1-1) + (1-1) + (1-1) + \cdots = 0 + 0 + 0 + \cdots $,其部分和恒为 0，因此和为 0。如果从第一项后开始成对分组，则得到：$1 + (-1+1) + (-1+1) + \cdots = 1 + 0 + 0 + \cdots $,其部分和恒为 1，因此和为 1，与前一种分组方式结果不同。

一般来说，将级数分组会得到一个新的级数，其部分和序列是原级数部分和序列的子序列。这意味着：如果原级数收敛，则分组后的新级数也收敛，并且极限相同——因为收敛数列的任意无限子序列也收敛到同一极限。但是，如果原级数发散，则分组后的级数不一定发散，就像上述格兰迪级数的例子。然而，如果分组后的级数发散，则必然说明原级数发散，因为这表明原级数部分和序列存在一个不收敛的子序列，而这在收敛的情况下是不可能的。

这一推理曾被奥雷姆用来证明调和级数发散[28]，并且是一般柯西凝聚判别法的基础。[29][30]
\subsubsection{重排}
在普通的有限求和中，由于加法的交换律，和式中的项可以任意重排，而不会改变求和的结果：$a_{0} + a_{1} + a_{2} = a_{0} + a_{2} + a_{1} = a_{2} + a_{1} + a_{0}$.类似地，在一个级数中，任何有限次重排都不会改变其部分和的极限，因此不会改变级数的和：对于任何有限的重排，总存在某一项之后，重排对后续项不再产生影响；这样，重排的影响仅限于有限项的求和，而有限和在重排下是不变的。

然而，就像分组一样，在一个无限级数中进行无限次重排，有时可能会改变部分和的极限。那些部分和序列收敛于某个值，但其项可以被重新排列成另一个级数，而该级数的部分和却收敛于不同值的情形，被称为条件收敛级数。而无论如何重排都收敛于同一个值的，则称为无条件收敛级数。

对于实数或复数项的级数：$a_{0} + a_{1} + a_{2} + \cdots$当且仅当其项的绝对值级数$|a_{0}| + |a_{1}| + |a_{2}| + \cdots$也收敛时，它是无条件收敛的。这一性质称为绝对收敛。否则，若一个实数或复数级数收敛但并非绝对收敛，它就是条件收敛的。

任何实数条件收敛级数都可以通过重新排列其项，使其和收敛到任意一个给定的实数，甚至可以使其发散。这一结论就是黎曼级数定理的内容。[31][32][33]

一个历史上重要的条件收敛的例子是交错调和级数：
$$
\sum_{n=1}^{\infty} \frac{(-1)^{n+1}}{n} 
= 1 - \frac{1}{2} + \frac{1}{3} - \frac{1}{4} + \frac{1}{5} - \cdots ,~
$$
它的和等于自然对数 $\ln 2$。而其项的绝对值之和则是**调和级数**：
$$
\sum_{n=1}^{\infty} \frac{1}{n} 
= 1 + \frac{1}{2} + \frac{1}{3} + \frac{1}{4} + \frac{1}{5} + \cdots ,~
$$
由于调和级数发散，[28] 所以交错调和级数是条件收敛的。例如，将交错调和级数的项重新排列，使得原级数中的每一个正项后面跟着两个负项（而不是一个），则得到：[34]
$$
1 - \tfrac{1}{2} - \tfrac{1}{4} + \tfrac{1}{3} - \tfrac{1}{6} - \tfrac{1}{8} + \tfrac{1}{5} - \tfrac{1}{10} - \tfrac{1}{12} + \cdots~
$$
$$
= \left(1 - \tfrac{1}{2}\right) - \tfrac{1}{4} + \left(\tfrac{1}{3} - \tfrac{1}{6}\right) - \tfrac{1}{8} + \left(\tfrac{1}{5} - \tfrac{1}{10}\right) - \tfrac{1}{12} + \cdots~
$$
$$
= \tfrac{1}{2} - \tfrac{1}{4} + \tfrac{1}{6} - \tfrac{1}{8} + \tfrac{1}{10} - \tfrac{1}{12} + \cdots~
$$
$$
= \tfrac{1}{2}\left(1 - \tfrac{1}{2} + \tfrac{1}{3} - \tfrac{1}{4} + \tfrac{1}{5} - \tfrac{1}{6} + \cdots\right),~
$$
这正好是原交错调和级数的 $\tfrac{1}{2}$，因此它的和是 $\tfrac{1}{2}\ln 2$。根据黎曼级数定理，通过重新排列交错调和级数的项，可以使其和变为任意一个实数，甚至使其发散。
\subsection{运算}
\subsubsection{级数的加法}
两个级数$a_{0} + a_{1} + a_{2} + \cdots \quad \text{和} \quad 
b_{0} + b_{1} + b_{2} + \cdots$的加法由逐项相加给出[13][35][36][37]：
$(a_{0}+b_{0}) + (a_{1}+b_{1}) + (a_{2}+b_{2}) + \cdots $,或者用求和符号表示：
$$
\sum_{k=0}^{\infty} a_{k} + \sum_{k=0}^{\infty} b_{k} 
= \sum_{k=0}^{\infty} (a_{k}+b_{k}).~
$$
若用 $s_{a,n}$、$s_{b,n}$ 分别表示被相加的两个级数的部分和，用 $s_{a+b,n}$ 表示结果级数的部分和，那么根据定义可得：$s_{a+b,n} = s_{a,n} + s_{b,n}$.因此，结果级数的和，即其部分和序列的极限，满足：
$$
\lim_{n \to \infty} s_{a+b,n} 
= \lim_{n \to \infty} (s_{a,n} + s_{b,n}) 
= \lim_{n \to \infty} s_{a,n} + \lim_{n \to \infty} s_{b,n},~
$$
在极限存在的情况下。由此可见：如果两个级数都是可和的，那么它们的和也是可和的；结果级数的和等于两个原级数的和之和。两个发散级数的加法有时可能得到一个收敛级数：例如，将一个发散级数与其各项乘以 $-1$ 的级数相加，会得到一个全为零的级数，它收敛到 0。然而，如果两个级数中一个收敛、一个发散，那么它们的和必然发散。[35]

对于实数级数或复数级数，级数加法是结合的、交换的和可逆的。因此，级数加法赋予实数或复数收敛级数集合以阿贝尔群的结构，同时也赋予所有实数或复数级数的集合（不论其收敛性如何）以阿贝尔群的结构。
\subsubsection{数乘}
一个级数$a_{0} + a_{1} + a_{2} + \cdots$与一个常数 $c$（在此情境下称为标量）的积，由逐项相乘给出[35]：$ca_{0} + ca_{1} + ca_{2} + \cdots $,或者用求和符号表示：
$$
c \sum_{k=0}^{\infty} a_{k} \;=\; \sum_{k=0}^{\infty} ca_{k}.~
$$
若用$s_{a,n}$表示原级数的部分和，用$s_{ca,n}$表示与标量$c$相乘后的级数的部分和，那么根据定义有：$s_{ca,n} = c \, s_{a,n}, \quad \text{对所有 } n$.因此也有：$\lim_{n \to \infty} s_{ca,n} = c \, \lim_{n \to \infty} s_{a,n}$,在极限存在的情况下。由此可见：如果一个级数是可和的，那么它的任意非零标量倍数也是可和的；反之，如果一个级数是发散的，那么它的任意非零标量倍数也是发散的。

实数和复数上的数乘运算是结合的、交换的、可逆的，并且对级数加法满足分配律。

综上所述，级数的加法与数乘运算使得收敛级数集合以及实数级数集合构成了一个实向量空间。同理，对于复数级数与收敛复数级数，则得到复向量空间。这些向量空间都是无限维的。
\subsubsection{级数的乘法}
两个级数$a_{0} + a_{1} + a_{2} + \cdots \quad \text{与} \quad b_{0} + b_{1} + b_{2} + \cdots$相乘生成第三个级数$c_{0} + c_{1} + c_{2} + \cdots$,这种乘法称为柯西乘积[12][13][14][36][38]。用求和符号表示为：
$$
\Biggl(\sum_{k=0}^{\infty} a_{k}\Biggr) \cdot \Biggl(\sum_{k=0}^{\infty} b_{k}\Biggr) 
= \sum_{k=0}^{\infty} c_{k} 
= \sum_{k=0}^{\infty} \sum_{j=0}^{k} a_{j} b_{k-j},~
$$
其中每个 $c_{k}$ 由下式给出：$c_{k} = \sum_{j=0}^{k} a_{j} b_{k-j} = a_{0} b_{k} + a_{1} b_{k-1} + \cdots + a_{k-1} b_{1} + a_{k} b_{0}.$在这里，结果级数$c_{0} + c_{1} + c_{2} + \cdots$的部分和是否收敛，并不像加法那样容易判定。然而，如果两个级数$a_{0} + a_{1} + a_{2} + \cdots \quad \text{与} \quad b_{0} + b_{1} + b_{2} + \cdots$都是绝对收敛的，那么它们的乘积级数也绝对收敛，其和等于原两个级数的和的乘积[13][36][39]：
$$
\lim_{n \to \infty} s_{c,n} 
= \Bigl(\lim_{n \to \infty} s_{a,n}\Bigr) \cdot \Bigl(\lim_{n \to \infty} s_{b,n}\Bigr).~
$$
对于实数或复数的绝对收敛级数，级数乘法是结合的、交换的，并且对级数加法满足分配律。因此，结合级数加法，级数乘法使得实数或复数的绝对收敛级数集合具有交换环的结构；再结合数乘，它们进一步构成交换代数的结构。同样，这些运算也使得所有实数或复数级数的集合具有结合代数的结构。
